\section{Билет 9. Квазиньютоновские методы}

\ticketq{Квазиньютоновские методы.}

\paragraph{Идея.}
Метод Ньютона требует на каждой итерации вычислять $H(x_k) = \nabla^2 f(x_k)$ и
решать $H_k p_k = -g_k$ — дорого. Квазиньютоновские методы гессиан не
вычисляют: вместо $H^{-1}(x_k)$ берётся приближение $A_k$,
\[
  p_k = -A_k \nabla f(x_k) ,
  \qquad
  x_{k+1} = x_k + \alpha_k p_k ,
  \qquad
  A_k \xrightarrow[k\to\infty]{} H^{-1}(x_*) ,
\]
причём $A_{k+1} = A_k + \Delta A_k$ и $A_1 = E$ (первая итерация —
градиентный шаг, информации о кривизне ещё нет). Методы занимают промежуточное
положение: используют информацию первого порядка, но имитируют поведение
второго.

\paragraph{Разности и секущее условие.}
\[
  s_k = x_{k+1} - x_k , \qquad
  y_k = \nabla f(x_{k+1}) - \nabla f(x_k) .
\]
Для точного гессиана $y_k \approx H_{k+1}s_k$, поэтому от приближения обратной
матрицы требуют \textbf{секущее (квазиньютоновское) условие}
\[
  A_{k+1}\,y_k = s_k ,
\]
то есть чтобы $A_{k+1}$ действовала на уже известном направлении так же, как
настоящая $H^{-1}$. \textbf{Условие кривизны:} $s_k\T y_k > 0$.

\paragraph{Теорема (направление спуска).}
\textit{$A_k \succ 0$, $g_k \ne 0$ $\Rightarrow$ $p_k = -A_k g_k$ —
направление спуска.}

\textsc{Док-во.} $g_k\T p_k = -g_k\T A_k g_k < 0$, так как $A_k \succ 0$ и
$g_k \ne 0$. $\blacksquare$

Поэтому положительная определённость $A_k$ — то, что требуется сохранять при
обновлении.

\paragraph{Формула ДФП (Дэвидон--Флетчер--Пауэлл).}
\[
  A_{k+1} = A_k + \frac{s_k s_k\T}{s_k\T y_k}
            - \frac{A_k y_k y_k\T A_k}{y_k\T A_k y_k} .
\]

\textbf{Теорема (секущее условие).} \textit{$A_{k+1}y_k = s_k$.}

\textsc{Док-во.} Умножим формулу справа на $y_k$:
\[
  A_{k+1}y_k = A_k y_k
    + \frac{s_k\,(s_k\T y_k)}{s_k\T y_k}
    - \frac{A_k y_k\,(y_k\T A_k y_k)}{y_k\T A_k y_k}
    = A_k y_k + s_k - A_k y_k = s_k ,
\]
так как $s_k\T y_k$ и $y_k\T A_k y_k$ — \emph{числа}. $\blacksquare$

\textbf{Теорема (положительная определённость).}
\textit{$A_k \succ 0$, $s_k\T y_k > 0$ $\Rightarrow$ $A_{k+1} \succ 0$.}

\textsc{Док-во.} Для $z \ne 0$
\[
  z\T A_{k+1} z = \underbrace{z\T A_k z
    - \frac{(z\T A_k y_k)^2}{y_k\T A_k y_k}}_{\ \ge\ 0}
    + \underbrace{\frac{(z\T s_k)^2}{s_k\T y_k}}_{\ \ge\ 0} \ \ge\ 0 ,
\]
первое слагаемое неотрицательно по Коши--Буняковскому в метрике
$\langle u,v\rangle_{A_k} = u\T A_k v$, второе — из $s_k\T y_k > 0$.
Равенство нулю требует обнуления обоих: из равенства в Коши следует
$z = c\,y_k$, а из $(z\T s_k)^2 = 0$ тогда $c\,y_k\T s_k = 0$. Но
$y_k\T s_k > 0$, значит $c = 0$ и $z = 0$. Поэтому при $z \ne 0$ строго
$z\T A_{k+1}z > 0$. $\blacksquare$

\paragraph{Формула БФГШ (Бройден--Флетчер--Голдфарб--Шенно).}
При $\rho_k = \dfrac{1}{y_k\T s_k}$
\[
  A_{k+1} = \big(E - \rho_k s_k y_k\T\big) A_k
            \big(E - \rho_k y_k s_k\T\big) + \rho_k s_k s_k\T .
\]

\textbf{Теорема (секущее условие).} \textit{$A_{k+1}y_k = s_k$.}

\textsc{Док-во.} $\rho_k(s_k\T y_k) = 1$ $\Rightarrow$
$(E - \rho_k y_k s_k\T)y_k = y_k - \rho_k y_k (s_k\T y_k) = 0$ $\Rightarrow$
первое слагаемое обнуляется. Остаётся
$\rho_k s_k s_k\T y_k = \frac{1}{y_k\T s_k}s_k (s_k\T y_k) = s_k$.
$\blacksquare$

\textbf{Теорема (положительная определённость).}
\textit{$A_k \succ 0$, $y_k\T s_k > 0$ $\Rightarrow$ $A_{k+1} \succ 0$.}

\textsc{Док-во.} Обозначим $w = (E - \rho_k y_k s_k\T)z$, тогда
\[
  z\T A_{k+1} z = w\T A_k w + \rho_k (z\T s_k)^2 \ \ge\ 0 .
\]
Равенство $\Rightarrow$ $w\T A_k w = 0$ и $(z\T s_k)^2 = 0$. Из
$A_k \succ 0$ следует $w = 0$; из $s_k\T z = 0$ получаем
$w = z - \rho_k y_k(s_k\T z) = z$ $\Rightarrow$ $z = 0$. $\blacksquare$

\paragraph{Алгоритм.}
\begin{csteps}
  \item Задать $x_0$, $\eps > 0$, $A_0 = E$, $k = 0$.
  \item $g_k = \nabla f(x_k)$; если $\norm{g_k} < \eps$ — конец,
        $x_* \approx x_k$.
  \item $p_k = -A_k g_k$.
  \item $\alpha_k = \argmin_{\alpha>0} f(x_k + \alpha p_k)$ (либо условия
        Армихо/Вульфа); $x_{k+1} = x_k + \alpha_k p_k$.
  \item $s_k = x_{k+1} - x_k$, $y_k = \nabla f(x_{k+1}) - \nabla f(x_k)$.
  \item Если $k$ принадлежит заданному множеству номеров обновления $\Pi_0$ —
        \textbf{рестарт}: $A_{k+1} = E$ (накопленная матрица сбрасывается,
        обычно раз в $n$ итераций). Иначе при $s_k\T y_k > 0$ обновить $A_k$
        по ДФП или БФГШ; при $s_k\T y_k \le 0$ — тоже сброс $A_{k+1} = E$.
        $k := k+1$, к~шагу~2.
\end{csteps}

\paragraph{Свойства и ограничения.}
Условие кривизны $s_k\T y_k > 0$ выполняется автоматически для сильно выпуклой
$f$ при точном одномерном поиске; в общем случае его обеспечивают условиями
Вульфа. БФГШ устойчивее ДФП: обновление записано как сопряжение
$W A_k W\T$ плюс слагаемое ранга один, и положительная определённость не
теряется от накопления ошибок округления, тогда как в ДФП она держится на
разности двух слагаемых. Для квадратичной функции $n$ переменных при точном
одномерном поиске метод сходится не более чем за $n$ шагов, и $A_n = H^{-1}$.

Ограничения: нужно хранить и обновлять матрицу $n \times n$ (при большом $n$
используют L-BFGS); чувствительность к точности одномерного поиска; для
невыпуклых функций — только локальный минимум.

\seefull{гл.~2, \S23 (численный пример ДФП там же)}
